\section{Conclusion and Future Work}

We presented \MLIR, a flexible and extensible infrastructure for compiler
construction. This paper described \MLIR's concrete
design, demonstrated its applicability to a range of
important domains, and described a number of original research and engineering
implications.

Looking ahead, we are eager to see how established compiler communities
(e.g.\ the Clang C and \Cpp compiler) as well domain experts can benefit from the
introduction of higher level, language-specific IRs. We are also interested to
see if \MLIR enables new approaches to teaching the art of compiler and IR
design, and hope to see entirely new areas of research catalyzed or accelerated
by this infrastructure.

\paragraph{Future directions}

There are multiple future directions being pursued for \MLIR.
In the ML and HPC space, these include inferring efficient Op implementations from reference rank-polymorphic specifications with symbolic shapes. It also involves enabling a wider range of data structures (sparse, graphs) and program transformations, bringing together symbolic reasoning, such as automatic differentiation and algorithmic simplification, with more conventional data flow and control flow-based optimizations. Beyond ML and HPC, one may consider \MLIR's applicability to other related domains, such as secure compilation, safety-critical systems, data analytics and graph processing, relational query optimization, etc.

Returning to the world of general-purpose languages, an obvious missing front-end
is a C++ mid-level
representation derived from Clang. Say, a ``CIL'' similar in spirit to Swift's SIL and Rust's MIR, that would facilitate the optimization of
common C++ idioms that currently need to be reconstructed from lowered code (e.g.,
treating \lstinline{std::vector} as an array rather than pointer manipulation). Supporting garbage-collected languages, higher-order and polymorphic type systems with type inference in \MLIR are open challenges as well.

Exploring parallelism and concurrency constructs in LLVM has been difficult,
primarily as the changes required are invasive and not easily layered (e.g.,
injecting metadata and inspecting all passes to ensure metadata is propagated
while losing optimization opportunities as the level of abstraction is too low).
With \MLIR, parallel constructs
can be first-class operations, using regions and parallel idiom-specific verification. This would support
higher-level transformations before lowering to, e.g., LLVM where regular
transformations can be performed on already lowered code.

Beyond debugging and testing, textual form of IR is
also useful for education. Additional tooling to show the interaction of
optimizations in high-performance compilation could demystify compilers for new
students. IR design is an integral part of developing a new compiler or
optimization framework, but many undergraduate compiler curricula do not
cover IR design. \MLIR provides opportunities for new approaches to such 
lessons that could be explored.
